\documentclass[10pt,a4paper]{amsart}
\usepackage{amsmath,amssymb,amsthm,graphicx}
% Retained minimal preamble for the Stacks Project chapter varieties.tex.
% Verbatim excerpt of the chapter preamble: Xy-pic setup (native preamble lines 14-18), the native theorem
% declarations with the mathematical macros (native preamble lines 154-217), and the project own no-op
% definitions of the reference/slogan/history display environments (native wrapper.tex lines 1-3) that the
% counted statement uses for its slogan. Copied mechanically; nothing added, removed or redefined.
% For commutative diagrams we use Xy-pic
\usepackage[all]{xy}

% We use 2cell for 2-commutative diagrams.
\xyoption{2cell}
% Theorem environments.
%
\theoremstyle{plain}
\newtheorem{theorem}[subsection]{Theorem}
\newtheorem{proposition}[subsection]{Proposition}
\newtheorem{lemma}[subsection]{Lemma}

\theoremstyle{definition}
\newtheorem{definition}[subsection]{Definition}
\newtheorem{example}[subsection]{Example}
\newtheorem{exercise}[subsection]{Exercise}
\newtheorem{situation}[subsection]{Situation}

\theoremstyle{remark}
\newtheorem{remark}[subsection]{Remark}
\newtheorem{remarks}[subsection]{Remarks}

\numberwithin{equation}{subsection}

% Macros
%
\def\lim{\mathop{\mathrm{lim}}\nolimits}
\def\colim{\mathop{\mathrm{colim}}\nolimits}
\def\Spec{\mathop{\mathrm{Spec}}}
\def\Hom{\mathop{\mathrm{Hom}}\nolimits}
\def\Ext{\mathop{\mathrm{Ext}}\nolimits}
\def\SheafHom{\mathop{\mathcal{H}\!\mathit{om}}\nolimits}
\def\SheafExt{\mathop{\mathcal{E}\!\mathit{xt}}\nolimits}
\def\Sch{\mathit{Sch}}
\def\Mor{\mathop{\mathrm{Mor}}\nolimits}
\def\Ob{\mathop{\mathrm{Ob}}\nolimits}
\def\Sh{\mathop{\mathit{Sh}}\nolimits}
\def\NL{\mathop{N\!L}\nolimits}
\def\CH{\mathop{\mathrm{CH}}\nolimits}
\def\proetale{{pro\text{-}\acute{e}tale}}
\def\etale{{\acute{e}tale}}
\def\QCoh{\mathit{QCoh}}
\def\Ker{\mathop{\mathrm{Ker}}}
\def\Im{\mathop{\mathrm{Im}}}
\def\Coker{\mathop{\mathrm{Coker}}}
\def\Coim{\mathop{\mathrm{Coim}}}

% Boxtimes
%
\DeclareMathSymbol{\boxtimes}{\mathbin}{AMSa}{"02}

%
% Macros for moduli stacks/spaces
%
\def\QCohstack{\mathcal{QC}\!\mathit{oh}}
\def\Cohstack{\mathcal{C}\!\mathit{oh}}
\def\Spacesstack{\mathcal{S}\!\mathit{paces}}
\def\Quotfunctor{\mathrm{Quot}}
\def\Hilbfunctor{\mathrm{Hilb}}
\def\Curvesstack{\mathcal{C}\!\mathit{urves}}
\def\Polarizedstack{\mathcal{P}\!\mathit{olarized}}
\def\Complexesstack{\mathcal{C}\!\mathit{omplexes}}
% \Pic is the operator that assigns to X its picard group, usage \Pic(X)
% \Picardstack_{X/B} denotes the Picard stack of X over B
% \Picardfunctor_{X/B} denotes the Picard functor of X over B
\def\Pic{\mathop{\mathrm{Pic}}\nolimits}
\def\Picardstack{\mathcal{P}\!\mathit{ic}}
\def\Picardfunctor{\mathrm{Pic}}
\def\Deformationcategory{\mathcal{D}\!\mathit{ef}}
\newenvironment{reference}{}{}
\newenvironment{slogan}{}{}
\newenvironment{history}{}{}
\usepackage{cleveref}
\crefname{originallemma}{original lemma}{original lemmas}
\Crefname{originallemma}{Original lemma}{Original lemmas}
\makeatletter
\newlabel{algebra-lemma-completion-generalities}{{0315}{source}{}{sourcepointer-1}{}}
\newlabel{algebra-lemma-completion-generalities@cref}{{[originallemma][0][]0315}{[1][0][]source}{}{}{}}
\makeatother
\crefname{originallemma}{original lemma}{original lemmas}
\Crefname{originallemma}{Original lemma}{Original lemmas}
\makeatletter
\newlabel{algebra-lemma-hathat-finitely-generated}{{05GG}{source}{}{sourcepointer-2}{}}
\newlabel{algebra-lemma-hathat-finitely-generated@cref}{{[originallemma][0][]05GG}{[1][0][]source}{}{}{}}
\makeatother
\crefname{originaltheorem}{original theorem}{original theorems}
\Crefname{originaltheorem}{Original theorem}{Original theorems}
\makeatletter
\newlabel{algebra-theorem-cohen-structure-theorem}{{032A}{source}{}{sourcepointer-3}{}}
\newlabel{algebra-theorem-cohen-structure-theorem@cref}{{[originaltheorem][0][]032A}{[1][0][]source}{}{}{}}
\makeatother
\crefname{originallemma}{original lemma}{original lemmas}
\Crefname{originallemma}{Original lemma}{Original lemmas}
\makeatletter
\newlabel{more-algebra-lemma-lift-residue-field}{{0C34}{source}{}{sourcepointer-4}{}}
\newlabel{more-algebra-lemma-lift-residue-field@cref}{{[originallemma][0][]0C34}{[1][0][]source}{}{}{}}
\makeatother
\title{Presentation of a complete local ring of a locally Noetherian scheme as a quotient of a power series ring, and when it is a k-algebra isomorphism}
\author{The Stacks Project Authors}
\date{stacks-project:varieties.tex:local-snapshot-20261011; GNU FDL 1.2 or later}
\begin{document}
\maketitle
\noindent\textit{Editorial note.} All delivered mathematical text is the author's own, copied from the pinned Stacks Project chapter file varieties.tex; no statement or proof step was rewritten, reordered or removed.\par\medskip
\noindent\textit{Editorial note.} Two of the author's footnotes belong to the counted proof span and are reproduced verbatim: the one explaining the mixed characteristic form of the Cohen structure theorem, and the aside pointing to the computation modulo y\textasciicircum{}2.\par\medskip
\noindent\textit{Editorial note.} The counted spans carry 7 navigation stubs and no bibliography citation.\par\medskip
\noindent\textit{Editorial note.} The preamble is the chapter's retained minimal preamble (sources/stacks-varieties/retained-preamble.tex lines 1-77), a verbatim excerpt of the native preamble.\par\medskip
\noindent\textit{Editorial note.} The label lemma-complete-local-ring-structure-as-algebra is defined only in varieties.tex, so no cross-chapter restatement had to be weighed for this card.\par\medskip
\noindent\textit{Allowed standard prerequisites (editorial).}\begin{itemize}
\item Cohen structure theorem for complete local rings. Algebra, Theorem algebra-theorem-cohen-structure-theorem, invoked at proof lines 4060-4067. Supplies the existence of the presentation by a quotient of a power series ring.
\item Lifting the residue field along a separable extension. More on Algebra, Lemma more-algebra-lemma-lift-residue-field, invoked at proof lines 4081-4084. Supplies the k-algebra homomorphism kappa to the completion when kappa/k is separable.
\item Completeness of a finitely generated module and the surjectivity criterion for adically complete maps. Algebra, Lemmas algebra-lemma-hathat-finitely-generated and algebra-lemma-completion-generalities, invoked at proof lines 4091-4096. Supplies that the constructed map of complete rings is surjective once it is surjective modulo the augmentation ideal.
\end{itemize}\medskip
\ifdefined\hypertarget\hypertarget{sourcepointer-1}{}\fi
\noindent\textit{Source reference.} Original lemma 0315: sources/stacks-chow/tags.txt, source lines 3972--3972. Pointer outside the counted span, invoked once for context; the load-bearing argument of the counted statement is the author's proof reproduced in full in the proof span.\par\medskip
\ifdefined\hypertarget\hypertarget{sourcepointer-2}{}\fi
\noindent\textit{Source reference.} Original lemma 05GG: sources/stacks-chow/tags.txt, source lines 7060--7060. Pointer outside the counted span, invoked once for context; the load-bearing argument of the counted statement is the author's proof reproduced in full in the proof span.\par\medskip
\ifdefined\hypertarget\hypertarget{sourcepointer-3}{}\fi
\noindent\textit{Source reference.} Original theorem 032A: sources/stacks-chow/tags.txt, source lines 4013--4013. Pointer outside the counted span, invoked once for context; the load-bearing argument of the counted statement is the author's proof reproduced in full in the proof span.\par\medskip
\ifdefined\hypertarget\hypertarget{sourcepointer-4}{}\fi
\noindent\textit{Source reference.} Original lemma 0C34: sources/stacks-chow/tags.txt, source lines 15179--15179. Pointer outside the counted span, invoked once for context; the load-bearing argument of the counted statement is the author's proof reproduced in full in the proof span.\par\medskip
\begin{lemma}
\label{lemma-complete-local-ring-structure-as-algebra}
Let $k$ be a field. Let $X$ be a locally Noetherian scheme over $k$.
Let $x \in X$ be a point with residue field $\kappa$.
There is an isomorphism
\begin{equation}
\label{equation-complete-local-ring}
\kappa[[x_1, \ldots, x_n]]/I \longrightarrow \mathcal{O}_{X, x}^\wedge
\end{equation}
inducing the identity on residue fields.
In general we cannot choose (\ref{equation-complete-local-ring})
to be a $k$-algebra isomorphism. However, if the extension $\kappa/k$
is separable, then we can choose
(\ref{equation-complete-local-ring}) to be an isomorphism of $k$-algebras.
\end{lemma}

\begin{proof}
The existence of the isomorphism is an immediate consequence of
the Cohen structure theorem\footnote{Note that if $\kappa$ has
characteristic $p$, then the theorem just says we get a surjection
$\Lambda[[x_1, \ldots, x_n]] \to \mathcal{O}_{X, x}^\wedge$ where
$\Lambda$ is a Cohen ring for $\kappa$. But of course in this
case the map factors through $\Lambda/p\Lambda[[x_1, \ldots, x_n]]$
and $\Lambda/p\Lambda = \kappa$.}
(Algebra, Theorem \ref{algebra-theorem-cohen-structure-theorem}).

\medskip\noindent
Let $p$ be an odd prime number, let $k = \mathbf{F}_p(t)$, and
$A = k[x, y]/(y^2 + x^p - t)$. Then the completion $A^\wedge$ of $A$
in the maximal ideal
$\mathfrak m = (y)$ is isomorphic to $k(t^{1/p})[[z]]$ as a ring
but not as a $k$-algebra. The reason is that $A^\wedge$
does not contain an element whose $p$th power is $t$ (as the reader
can see by computing modulo $y^2$).
This also shows that any isomorphism (\ref{equation-complete-local-ring})
cannot be a $k$-algebra isomorphism.

\medskip\noindent
If $\kappa/k$ is separable, then there is a $k$-algebra
homomorphism $\kappa \to \mathcal{O}_{X, x}^\wedge$
inducing the identity on residue fields by
More on Algebra, Lemma \ref{more-algebra-lemma-lift-residue-field}.
Let $f_1, \ldots, f_n \in \mathfrak m_x$ be generators.
Consider the map
$$
\kappa[[x_1, \ldots, x_n]] \longrightarrow \mathcal{O}_{X, x}^\wedge,\quad
x_i \longmapsto f_i
$$
Since both sides are $(x_1, \ldots, x_n)$-adically complete
(the right hand side by
Algebra, Lemmas \ref{algebra-lemma-hathat-finitely-generated})
this map is surjective by
Algebra, Lemma \ref{algebra-lemma-completion-generalities}
as it is surjective modulo $(x_1, \ldots, x_n)$ by
construction.
\end{proof}

\end{document}
